\begin{textsample}{POS Dim 4 – 2013\_c6 – Score 8.00 – 2013\_c6/enade\_2013\_zoot\_9cdde507ba.txt}  \label{ex:f4_pos_047}
CAPÍTULO 3 ANÁLISE TÉCNICA DA PROVA Este capítulo tem por objetivo apresentar o desempenho dos estudantes concluintes de Zootecnia no ENADE/2013.
Para isso, foram calculadas as estatísticas básicas da prova em seu todo ( seção 3.1.1 ), bem como as estatísticas dos componentes relacionadas ao Componente de Formação Geral ( seção 3.1.2 ) e ao de Conhecimento Específico da Área ( seção 3.1.3 ).
Dadas as suas características, foram analisadas em separado, as questões objetivas ( seção 3.2 ) e as questões discursivas ( seção 3.3 ).
Para as questões objetivas, foram disponibilizados os índices de facilidade e o de discriminação bisserial, também em separado para os Componentes de Formação Geral ( seção 3.2.1 ) e de Conhecimento Específico ( seção 3.2.2 ).
De cada componente, uma das questões foi escolhida para exemplificar a análise gráfica relacionando as alternativas escolhidas pelos estudantes ( inclusive o gabarito ) com o número de acertos no componente.
O Anexo I apresenta a íntegra da análise gráfica para todas as questões objetivas.
Para cada uma das questões discursivas, os conteúdos dos tipos mais comuns de respostas dos alunos são apresentados e comparados com o padrão de resposta esperado ( ver Anexo VIII com o padrão de respostas ).
Tomando como base as duas questões discursivas do Componente de Formação Geral, a seção 3.3.1.6 apresenta comentários sobre a correção das respostas com respeito à Língua Portuguesa.
Nas tabelas, são apresentados o tamanho da população inscrita e de presentes, e as seguintes estatísticas das notas11 : média do desempenho na prova, erro padrão da média, desvio padrão, nota mínima, mediana e nota máxima.
Tais estatísticas contemplam o total de estudantes concluintes da Área de Zootecnia inscritos e presentes à prova do ENADE/2013, tendo -se em vista as seguintes agregações : ( a ) as Grandes Regiões e o país como um todo ; ( b ) a Categoria Administrativa ; e ( c ) a Organização Acadêmica.
Em relação aos gráficos de distribuição de notas, o intervalo considerado foi de 10 unidades, aberto à esquerda e fechado à direita, com exceção do primeiro intervalo, [ 0 ; 10 ], fechado em ambos os extremos.
Para os gráficos de distribuição das notas das questões discursivas, foram consideradas mais duas categorias : questão em branco e nota zero. 11 Essas estatísticas e outras estão definidas no Capítulo 1. 3.1 ESTATÍSTICAS BÁSICAS DA PROVA 3.1.1 Estatísticas Básicas Gerais A Tabela 3.1 apresenta as estatísticas básicas da prova por grande Região.
A população total de inscritos foi de 2.349.
Destes, 2.109 estiveram presentes, sendo 10,2 \% o índice de não comparecimento.
A Região de maior abstenção foi a Sudeste ( 14,9 \% ) e a de menor abstenção foi a Sul ( 4,3 \% ).
A média das notas da prova como um todo ( nas seções seguintes serão analisados os Componentes de Formação Geral e de Conhecimento Específico ) foi 46,8, sendo que os alunos da região Norte obtiveram a média mais baixa ( 41,1 ) e os da região Sudeste obtiveram a média mais alta ( 51,5 ).
As demais médias foram : 42,9 na região Nordeste, 50,1 na região Sul e 45,3 na região Centro-Oeste.
O desvio padrão para o Brasil como um todo foi 12,8, sendo o maior desvio padrão encontrado na região Nordeste ( 12,7 ) e o menor nas regiões Sudeste e Sul ( 11,8 ), indicando uma menor dispersão das notas destas últimas regiões.
As regiões que obtiveram a maior nota máxima foram a Nordeste e a Sudeste ( 84,9 ), ao passo que a região que atingiu a menor nota máxima foi a Norte ( 71,6 ).
A mediana do Brasil como um todo foi 46,9, sendo a maior mediana obtida na região Sudeste ( 51,3 ) e a menor obtida na Norte ( 40,9 ).
A menor nota mínima foi obtida na região Nordeste ( 3,8 ), enquanto a maior nota mínima foi obtida na região Sudeste ( 16,7 ).
O comportamento das notas dos estudantes de todo o Brasil pode ser observado no Gráfico 3.1 que apresenta um histograma com a distribuição das mesmas.
Essa é uma distribuição unimodal com moda no intervalo ( 40 ; 50 ].
O coeficiente de assimetria da distribuição das notas é negativo ( -0,02 ), porém de pouca magnitude, e podemos considerar que a distribuição é aproximadamente simétrica.
As distribuições nas regiões Norte e Sudeste também apresentam assimetria negativa e pequena, concentração pouco maior do lado direito do histograma e mais espalhada do lado esquerdo.
Já as distribuições nas regiões Nordeste, Sul e Centro-Oeste apresentam assimetria positiva, concentração pouco maior do lado esquerdo do histograma e mais espalhada do lado direito.
Os Gráficos 3.2, 3.3 e 3.4 apresentam informações referentes à média da nota final dos participantes, desagregando os resultados de acordo com, respectivamente, as Grandes Regiões do país, a Categoria Administrativa e a Organização Acadêmica.
Os gráficos apresentam o valor da média das notas como uma barra e os extremos do intervalo de confiança de 95 \% como linhas verticais unidas por uma linha horizontal na forma da letra H maiúscula.
Considerando -se o gráfico de notas segundo Grande Região ( Gráfico 3.2 ), observa -se que existe diferença estatisticamente significativa ao nível de 95 \% entre a maior média, obtida na região Sudeste ( 51,5 ) e a menor, obtida na região Norte ( 41,1 ).
Levando -se em conta os estudantes por Categorias Administrativas ( Gráfico 3.3 ), observa -se que existe diferença estatisticamente significante entre as médias das notas das IES Públicas e Privadas.
Pode ser observado ainda que a média dos alunos de IES Públicas ( 47,7 ) é maior que a dos alunos de IES Privadas ( 40,4 ).
A diferença entre as médias das regiões Sudeste e Norte ( 10,4 ), a maior e a menor média, é superior a diferença entre IES Públicas e Privadas ( 7,3 ).
Tendo como foco o Gráfico 3.4, que apresenta as notas médias das provas segundo Organização Acadêmica, constata -se que a diferença da média das notas dos estudantes provenientes de Universidades é estatisticamente significativa ao nível de 95 \% da média das notas dos estudantes provenientes de Faculdades.
Já a diferença entre a média dos estudantes de Universidades e Centros Universitários, assim como a diferença entre a média dos estudantes de Centros Universitários e Faculdades, não é estatisticamente significativa.
A maior média foi obtida pelos estudantes de Universidades ( 47,5 ), e a menor, pelos de Faculdades ( 39,0 ). 3.1.2 Estatísticas Básicas no Componente de Formação Geral A Tabela 3.2 apresenta as estatísticas básicas em relação ao componente da prova que avalia a Formação Geral dos estudantes concluintes.
Os alunos de todo Brasil obtiveram desempenho médio de 47,2.
Quanto à variabilidade, o desvio padrão das notas dos estudantes do Brasil como um todo foi 14,5.
A maior média foi obtida na região Sudeste ( 50,6 ), e a menor, na região Norte ( 43,4 ).
As demais médias foram : 45,7 na região Nordeste, 48,3 na região Sul e 45,1 na região Centro-Oeste.
Já o maior desvio padrão foi obtido na região Norte ( 15,6 ) e o menor na região Sul ( 13,5 ).
Os demais desvios padrões foram : 15,0 na região Nordeste, 13,7 na região Sudeste e 13,9 na região Centro-Oeste.
A maior nota no Componente de Formação Geral da prova do ENADE foi obtida por pelo menos um aluno da região Sul ( 89,6 ) enquanto que a menor nota máxima foi obtida na região Norte ( 79,6 ).
Nas outras regiões as notas máximas foram : 88,2 na região Nordeste, 86,4 na região Sudeste e 83,9 na região Centro-Oeste.
A mediana do Brasil como um todo foi 47,8, sendo a menor mediana encontrada na região Norte ( 43,1 ) e a maior encontrada na região Sudeste ( 51,7 ).
A menor nota mínima nesta parte foi zero na região Nordeste, O Gráfico 3.5 propicia a avaliação do desempenho dos estudantes no Componente de Formação Geral a partir do histograma da distribuição das notas correspondentes.
A distribuição é unimodal, com moda nos intervalos ( 40 ; 50 ] e ( 50 ; 60 ], enquanto na prova como um todo a moda foi alcançada apenas no intervalo ( 40 ; 50 ].
Nota -se, ainda, que no Gráfico 3.5 as notas apresentam uma maior dispersão do que no Gráfico 3.1 ( distribuição das notas da prova ), confirmado pela comparação dos desvios padrões : 12,8 para a nota da prova como um todo e 14,5 para o Componente de Formação Geral.
Para o Componente de Formação Geral, o coeficiente de assimetria da distribuição das notas dos estudantes é negativo ( -0,21 ).
Por ser um coeficiente pequeno, a distribuição apresenta uma leve concentração à direita e cauda maior à esquerda, porém não podemos considerar que a distribuição é aproximadamente simétrica.
Em quase todas as Grandes Regiões os histogramas também possuem assimetria levemente negativa ( entre -0,24 para a Sul e -0,16 para a Centro-Oeste ), sendo a região Norte a exceção com coeficiente nulo. enquanto a maior nota mínima foi obtida na região Centro-Oeste ( 10,0 ).
As demais notas mínimas foram 7,5 tanto na região Norte quanto nas regiões Sudeste e Sul.
O Gráfico 3.5 propicia a avaliação do desempenho dos estudantes no Componente de Formação Geral a partir do histograma da distribuição das notas correspondentes.
A distribuição é unimodal, com moda nos intervalos ( 40 ; 50 ] e ( 50 ; 60 ], enquanto na prova como um todo a moda foi alcançada apenas no intervalo ( 40 ; 50 ].
Nota -se, ainda, que no Gráfico 3.5 as notas apresentam uma maior dispersão do que no Gráfico 3.1 ( distribuição das notas da prova ), confirmado pela comparação dos desvios padrões : 12,8 para a nota da prova como um todo e 14,5 para o Componente de Formação Geral.
Para o Componente de Formação Geral, o coeficiente de assimetria da distribuição das notas dos estudantes é negativo ( -0,21 ).
Por ser um coeficiente pequeno, a distribuição apresenta uma leve concentração à direita e cauda maior à esquerda, porém não podemos considerar que a distribuição é aproximadamente simétrica.
Em quase todas as Grandes Regiões os histogramas também possuem assimetria levemente negativa ( entre -0,24 para a Sul e -0,16 para a Centro-Oeste ), sendo a região Norte a exceção com coeficiente nulo.
Nos Gráficos 3.6, 3.7 e 3.8 são apresentadas as informações referentes ao desempenho dos concluintes no Componente de Formação Geral, em diferentes agregações : Grande Região do país, Categoria Administrativa e Organização Acadêmica.
Observa -se pelo Gráfico 3.6 que existe diferença estatisticamente significativa ao nível de 95 \% entre a maior média das notas na Componente de Formação Geral, obtida na região Sudeste ( 50,6 ) e a menor, obtida na região Norte ( 43,4 ).
Vemos que o intervalo de confiança mais largo é o da região Norte ; já o intervalo mais estreito é observado na região Sudeste.
Este fato está relacionado, também, com o tamanho da população envolvida, menor na região Norte do que na Sudeste.
No Gráfico 3.7, que representa as notas médias no Componente de Formação Geral segundo Categoria Administrativa do país, observa -se que existe diferença estatisticamente significativa entre as médias.
Assim como ocorreu para as notas da prova como um todo, para o Componente de Formação Geral, os concluintes de Zootecnia, das IES Públicas ( 47,9 ) obtiveram uma média maior do que os das IES Privadas ( 41,8 ).
Considerando -se o tipo de Organização Acadêmica, nota -se, no Gráfico 3.8, que a diferença da média das notas dos estudantes provenientes de Universidades é estatisticamente significativa ao nível de 95 \% da média das notas dos estudantes provenientes de Faculdades.
Já a diferença entre a média dos estudantes de Universidades e Centros Universitários, assim como a diferença entre a média dos estudantes de Centros Universitários e Faculdades, não é estatisticamente significativa.
Nas Universidades, a média ( 47,7 ) é mais elevada do que em Centros Universitários ( 45,4 ) e Faculdades ( 41,1 ). 3.1.3 Estatísticas Básicas do Componente de Conhecimento Específico A Tabela 3.3 apresenta as estatísticas básicas referentes ao Componente de Conhecimento Específico da área de Zootecnia.
A média do desempenho dos alunos do Brasil como um todo foi 46,7.
A maior média foi obtida na região Sudeste ( 51,8 ), e a menor, na região Norte ( 40,3 ).
As demais médias foram : 42,0 na região Nordeste, 50,7 na região Sul e 45,4 na região Centro-Oeste.
Quanto à variabilidade das notas, o desvio padrão do Brasil como um todo foi 14,8, sendo o maior desvio padrão observado na região Nordeste ( 14,6 ) e o menor na região Sudeste ( 13,8 ).
Os demais desvios padrões foram : 14,2 na região Norte, 13,9 na região Sul e 14,3 na região Centro-Oeste.
A mediana das notas dos estudantes de todo o Brasil foi 46,9.
A maior mediana ocorreu na região Sudeste ( 52,2 ) e a menor na região Norte ( 40,6 ).
As demais medianas foram 41,1 na região Nordeste, 50,4 na região Sul e 45,4 na região Centro-Oeste.
A maior nota máxima do Brasil como um todo foi 89,4, sendo obtida por pelo menos um aluno da região Nordeste.
As demais notas máximas foram : 76,8 na região Norte, 87,5 na região Sudeste, 87,6 na região Sul e 85,1 na região Centro-Oeste.
A menor nota mínima foi zero na Assim como os Gráficos 3.1 e 3.5, o Gráfico 3.9, apresentado a seguir, proporciona uma avaliação do desempenho de concluintes em relação ao Componente de Conhecimento Específico com um histograma da distribuição das notas correspondentes.
Esta também é uma distribuição unimodal, e o grupo modal é o ( 40 ; 50 ], um dos grupos pertencentes à moda para a Formação Geral.
O coeficiente de assimetria da distribuição das notas do Componente de Conhecimento Específico é positivo ( 0,03 ).
Nota -se pelo histograma ( Gráfico 3.9 ) que esta distribuição apresenta uma leve concentração à esquerda e cauda maior à direita.
Em quase todas as Grandes Regiões os histogramas também possuem assimetria positiva ( entre 0,08 para o Norte e 0,17 para o Centro-Oeste ), excetuando -se a região Sudeste ( -0,12 ) que possui assimetria negativa. região Nordeste.
As demais notas mínimas foram : 11,1 na região Norte, 11,8 na região Sudeste, 14,8 na região Sul e 9,9 na região Centro-Oeste.
Assim como os Gráficos 3.1 e 3.5, o Gráfico 3.9, apresentado a seguir, proporciona uma avaliação do desempenho de concluintes em relação ao Componente de Conhecimento Específico com um histograma da distribuição das notas correspondentes.
Esta também é uma distribuição unimodal, e o grupo modal é o ( 40 ; 50 ], um dos grupos pertencentes à moda para a Formação Geral.
O coeficiente de assimetria da distribuição das notas do Componente de Conhecimento Específico é positivo ( 0,03 ).
Nota -se pelo histograma ( Gráfico 3.9 ) que esta distribuição apresenta uma leve concentração à esquerda e cauda maior à direita.
Em quase todas as Grandes Regiões os histogramas também possuem assimetria positiva ( entre 0,08 para o Norte e 0,17 para o Centro-Oeste ), excetuando -se a região Sudeste ( -0,12 ) que possui assimetria negativa.
Os Gráficos 3.10, 3.11 e 3.12 apresentam, respectivamente, uma comparação dos resultados em relação à Grande Região do país, à Categoria Administrativa e à Organização Acadêmica, agora levando em conta o desempenho de alunos no Componente de Conhecimento Específico da prova.
Pelo Gráfico 3.10, observa -se que existe diferença estatisticamente significativa entre a maior e a menor média das notas no Componente de Conhecimento Específico, obtidas nas regiões Sudeste ( 51,8 ) e Norte ( 40,3 ) respectivamente.
Quanto à Categoria Administrativa ( Gráfico 3.11 ), observa -se um comportamento semelhante àquele da Parte de Formação Geral e à prova como um todo, ou seja, existe diferença estatisticamente significativa entre as médias das IES Públicas e Privadas, sendo que a maior média foi obtida por alunos de IES Públicas de ensino ( 47,6 ).
Quanto ao Gráfico 3.12, observa -se, mais uma vez, que a diferença da média das notas dos estudantes provenientes de Universidades é estatisticamente significativa ao nível de 95 \% da média das notas dos estudantes provenientes de Faculdades.
Já a diferença entre a média dos estudantes de Universidades e Centros Universitários, assim como a diferença entre a média dos estudantes de Centros Universitários e Faculdades, não é estatisticamente significativa.
Sendo que a média dos concluintes das Universidades ( 47,5 ) foi maior do que de Centros Universitários ( 41,3 ) e de Faculdades ( 38,2 ). 3.2 ANÁLISE DAS QUESTÕES OBJETIVAS 3.2.1 Componente de Formação Geral A Tabela 3.4 apresenta as estatísticas básicas relativas às oito questões objetivas do componente da prova que abrange a Formação Geral dos estudantes.
A média do Brasil foi 51,1.
A maior média foi encontrada na região Sudeste ( 54,3 ) e a menor na região Norte ( 47,1 ).
As demais médias foram 50,2 na região Nordeste, 53,0 na região Sul e 47,7 na região Centro-Oeste.
O desvio padrão do Brasil foi 17,6, sendo o maior desvio padrão encontrado na região Nordeste ( 18,1 ) e o menor na região Sul ( 16,8 ).
Os demais desvios foram : 17,9 na região Norte, 17,1 na região Sudeste e 17,3 na região Centro-Oeste.
A mediana ( 50,0 ), nas questões objetivas do Componente de Formação Geral, foram as mesmas em todas as regiões do Brasil, sem exceção.
As notas máximas ( 100,0 ) foram obtidas em quase todas as regiões, exceto nas regiões Norte e Centro-Oeste ( 87,5 ).
As notas mínimas ( 0,0 ) foram obtidas em quase todas as regiões, exceto nas regiões Norte e Sul ( 12,5 ).
A Tabela 3.5 apresenta o índice de facilidade e o índice de discriminação ( ponto bisserial ) para cada uma das questões objetivas do Componente de Formação Geral.
Quanto ao índice de facilidade, foram usadas as seguintes cores para diferenciar o nível de dificuldade da questão : • Azul para as questões classificadas com índice muito fácil ( índice \&gt ; =0,86 ), verde para as questões classificadas com índice fácil ( 0,61 a 0,85 ), amarelo para as questões classificadas com médio ( 0,41 a 0,60 ), vermelho para as questões classificadas com difícil ( 0,16 a 0,40 ) e roxo para as questões classificadas com muito difícil ( \&lt ; =0,15 ).
Já quanto ao índice de discriminação, foram usadas as seguintes cores para qualificar a questão : • As questões classificadas com índice fraco receberam a cor vermelha ( índice \&lt ; =0,19 ), as classificadas com médio receberam a cor amarelo ( 0,20 a 0,29 ), as classificadas com bom receberam a cor verde ( 0,30 a 0,39 ) e as classificadas com muito bom ( \&gt ; =0,40 ) receberam a cor azul.
As questões objetivas do Componente de Formação Geral, segundo o índice de facilidade, foram assim avaliadas : das oito questões, uma teve o índice de facilidade classificado como muito fácil e duas como fácil.
Três questões foram tidas como médio, por terem índice de acertos situado na faixa entre 0,41 e 0,60 ( de 41,0 \% a 60,0 \% de acertos ).
Duas questões foram consideradas de dificuldade difícil, situando -se no intervalo entre 0,16 e 0,40 do índice de facilidade, ou seja, houve entre 16,0 \% e 40,0 \% de acertos.
Como já comentado, para análise das questões objetivas relativas à Formação Geral segundo o poder de discriminação, utilizou -se, o índice de discriminação ponto bisserial.
Nesta análise, as questões foram assim avaliadas : quatro das oito questões apresentaram A Tabela 3.5 apresenta o índice de facilidade e o índice de discriminação ( ponto bisserial ) para cada uma das questões objetivas do Componente de Formação Geral.
Quanto ao índice de facilidade, foram usadas as seguintes cores para diferenciar o nível de dificuldade da questão :  Azul para as questões classificadas com índice muito fácil ( índice \&gt ; =0,86 ), verde para as questões classificadas com índice fácil ( 0,61 a 0,85 ), amarelo para as questões classificadas com médio ( 0,41 a 0,60 ), vermelho para as questões classificadas com difícil ( 0,16 a 0,40 ) e roxo para as questões classificadas com muito difícil ( \&lt ; =0,15 ).
Já quanto ao índice de discriminação, foram usadas as seguintes cores para qualificar a questão :  As questões classificadas com índice fraco receberam a cor vermelha ( índice \&lt ; =0,19 ), as classificadas com médio receberam a cor amarelo ( 0,20 a 0,29 ), as classificadas com bom receberam a cor verde ( 0,30 a 0,39 ) e as classificadas com muito bom ( \&gt ; =0,40 ) receberam a cor azul.
As questões objetivas do Componente de Formação Geral, segundo o índice de facilidade, foram assim avaliadas : das oito questões, uma teve o índice de facilidade classificado como muito fácil e duas como fácil.
Três questões foram tidas como médio, por terem índice de acertos situado na faixa entre 0,41 e 0,60 ( de 41,0 \% a 60,0 \% de acertos ).
Duas questões foram consideradas de dificuldade difícil, situando -se no intervalo entre 0,16 e 0,40 do índice de facilidade, ou seja, houve entre 16,0 \% e 40,0 \% de acertos.
Como já comentado, para análise das questões objetivas relativas à Formação Geral segundo o poder de discriminação, utilizou -se, o índice de discriminação ponto bisserial.
Nesta análise, as questões foram assim avaliadas : quatro das oito questões apresentaram O Gráfico 3.13, para exemplificar, analisa o comportamento da questão de número 3 de Formação Geral.
Trata -se da terceira questão mais fácil e a que obteve o maior índice de discriminação dessa parte da prova.
Neste gráfico, cada uma das cinco curvas representa o percentual de respostas em determinada alternativa da questão, em função da nota dos estudantes nesta parte da prova ( Formação Geral/Múltipla Escolha ), antes de possíveis eliminações pelo critério do ponto bisserial.
A curva em vermelho corresponde à alternativa E, a correta para esta questão.
Assim, observa -se que entre os estudantes com notas mais baixas, nessa parte do exame, a situação mais frequente foi a escolha de uma das alternativas incorretas : a alternativa A ( em azul ) ou D ( em roxo ).
Na medida em que a nota aumenta, indicando desempenho índices acima de 0,40 e, assim, foram classificadas com índice muito bom para esse grupo de alunos ; três questões tiveram índice de discriminação bom, entre 0,30 e 0,39 ; e uma questão teve índice de discriminação médio, entre 0,20 e 0,29.
Nenhuma questão foi classificada com índice fraco e, portanto, nenhuma questão foi eliminada pelo critério ponto bisserial.
O índice de facilidade variou de 0,23 a 0,96, e o de discriminação, de 0,25 a 0,48.
As questões com índice de discriminação muito bom, de números 3, 6, 7 e 8, figuraram entre as mais fáceis desse conjunto : uma ( questão 8 ) classificada na categoria fácil e outras três na categoria médio do índice de facilidade.
Em particular, a questão 3 apresentou maior poder discriminatório, com índice 0,48, e foi a quarta mais fácil, com uma proporção de 0,47 de acertos.
A questão de número 5 apresentou índice de facilidade 0,23, ou seja, um quantitativo de 23,0 \% de estudantes conseguiu resolvê -la, dentro do universo de participantes.
Além disso, seu índice de discriminação foi bom, com valor igual a 0,36.
O Gráfico 3.13, para exemplificar, analisa o comportamento da questão de número 3 de Formação Geral.
Trata -se da terceira questão mais fácil e a que obteve o maior índice de discriminação dessa parte da prova.
Neste gráfico, cada uma das cinco curvas representa o percentual de respostas em determinada alternativa da questão, em função da nota dos estudantes nesta parte da prova ( Formação Geral/Múltipla Escolha ), antes de possíveis eliminações pelo critério do ponto bisserial.
A curva em vermelho corresponde à alternativa E, a correta para esta questão.
Assim, observa -se que entre os estudantes com notas mais baixas, nessa parte do exame, a situação mais frequente foi a escolha de uma das alternativas incorretas : a alternativa A ( em azul ) ou D ( em roxo ).
Na medida em que a nota aumenta, indicando desempenho melhor nesta parte da prova, aumenta concomitantemente a proporção de estudantes que selecionaram a alternativa correta E, atingindo 100 \% para os alunos com 8 acertos.
Essa análise permite verificar como a questão discriminou os grupos de desempenho, justificando o alto índice obtido na questão.
Os gráficos relativos às demais questões de Formação Geral constam do Anexo I. 3.2.2 Componente de Conhecimento Específico A Tabela 3.6 apresenta as estatísticas básicas em relação às questões objetivas do Componente de Conhecimento Específico da prova, por Grande Região.
A média do Brasil deste componente foi de 49,6.
A menor média foi observada na região Norte ( 43,6 ) e a maior na região Sudeste ( 54,7 ).
O desvio padrão de todo o Brasil foi 15,7, sendo o menor desvio padrão encontrado na região Sudeste ( 14,6 ) e o maior na região Nordeste ( 15,7 ).
A mediana de todo o Brasil foi 47,8.
Na região Centro-Oeste, a mediana foi também 47,8.
A maior mediana foi encontrada na região Sudeste ( 56,5 ), enquanto a menor mediana foi encontrada na região Nordeste ( 43,5 ).
A maior nota máxima da prova ( 95,7 ) foi obtida, nas questões objetivas do Componente de Conhecimento Específico, por pelo menos um aluno na região Sul, enquanto a menor nota máxima foi obtida na região Norte ( 78,3 ).
A menor nota mínima foi zero na região Nordeste, enquanto a maior nota mínima foi obtida na região Sul ( 17,4 ).
A Tabela 3.7 apresenta o índice de facilidade e o índice de discriminação ( ponto bisserial ) para cada uma das questões objetivas do Componente de Conhecimento Específico.
Para facilitar a diferenciação das questões usaram -se as mesmas cores da Tabela 3.5 para as diferentes classificações dos índices de facilidade e de discriminação.
Dentre as questões objetivas da parte da prova relativa ao Componente de Conhecimento Específico, nenhuma foi anulada pela Comissão.
Desse modo, a classificação quanto ao índice de facilidade foi estabelecida com base nas 27 questões.
A partir dos índices obtidos, pode -se concluir que a maioria das 27 questões objetivas da prova foi considerada como média, 14 ao todo.
As 13 questões restantes foram classificadas quanto ao índice de facilidade da seguinte maneira : quatro como fácil e nove como difícil.
Já quanto aos índices de discriminação das questões objetivas do Componente de Conhecimento Específico da prova, tem -se como resultado a seguinte classificação : nove das 27 questões foram consideradas como bom, enquanto cinco delas tiveram índice de discriminação muito bom.
Assim, para 14 das 27 questões, os índices de discriminação foram bom ou muito bom.
Dentre as demais, nove delas foram classificadas como médio e outras quatro como fraco, sendo treze, por conseguinte, a quantidade de questões nos dois patamares mais baixos de discriminação.
Constata -se, assim, que a prova – no que se refere ao Componente de Conhecimento Específico – possuía capacidade mediana de discriminar entre aqueles que dominam ou não o conteúdo.
Dentre as questões que alcançaram os maiores índices de discriminação, as de números 9, 10, 26, 30 e 33 classificadas com índice muito bom, igual ou superior a 0,40, uma delas ( questão 10 ) foi classificada na categoria fácil e outras quatro como médio, quanto ao índice de facilidade.
A questão de número 29 foi a mais difícil dentre as 27 questões específicas, com baixo índice de facilidade, apenas 19,0 \% de acertos.
Essa questão apresentou poder discriminatório igualmente baixo, 0,16, o que comprova ter sido esta a mais difícil para os estudantes.
Destaca -se, também, a questão 24, com índice de facilidade 0,27, o que, em termos percentuais, corresponde a 27,0 \% de estudantes que responderam acertadamente, obtendo, ainda, 0,19 de índice de discriminação.
Tais questões foram, portanto, duas das mais difíceis da prova.
As questões 24 e 29 foram eliminadas do cômputo da nota final pelo critério do ponto bisserial.
Além destas duas, as demais questões com índice fraco de discriminação, questões 16 e 25 também não foram computadas.
A título de exemplo das análises do comportamento das questões objetivas, o Gráfico 3.14 analisa a questão 10 do Componente de Conhecimento Específico.
Esta foi uma das questões classificadas como fácil nesta prova, apresentando índice de facilidade 0,65, ou seja, 65,0 \% dos estudantes assinalaram acertadamente a opção C, correspondente ao gabarito.
Seu índice de discriminação foi igual a 0,46, classificado como muito bom.
Neste gráfico, cada uma das cinco curvas representa o percentual de respostas em determinada alternativa da questão 10, em função da nota dos estudantes nesta parte da prova, antes de possíveis eliminações de questões pelo critério do ponto bisserial.
A alternativa correta C, representada no gráfico pela curva em preto, foi escolhida em maiores proporções pelos alunos com desempenho melhor nesta parte da prova.
Já as alternativas incorretas, também denominadas distratores, foram selecionadas principalmente por aqueles com notas mais baixas.
A proporção de alunos que selecionou a resposta correta C aumenta gradativamente, chegando a atingir 100 \% para aqueles que acertaram pelo menos 20 questões, enquanto a proporção dos que escolheram alternativas incorretas decai, a partir da segunda nota não nula, como função da nota nesta parte da prova.
Os gráficos relativos às demais questões do Componente de Conhecimento Específico constam do Anexo I. 3.3 ANÁLISE DAS QUESTÕES DISCURSIVAS 3.3.1 Componente de Formação Geral As análises dos resultados de desempenho dos estudantes de Zootecnia nas duas questões discursivas relativas à Formação Geral encontram -se na Tabela 3.8 e no Gráfico 3.15.
Na Tabela 3.8 observa -se que as notas médias foram mais baixas nesse conjunto de questões do que no das objetivas.
Os estudantes de todo o Brasil obtiveram, em Formação Geral, média 51,1 nas questões objetivas e 41,3 nas questões discursivas.
A mediana também confirma o pior desempenho dos alunos de todo o Brasil nas questões discursivas do Componente de Formação Geral.
Enquanto essa estatística foi de 46,0 para questões discursivas, para as questões objetivas a mediana foi de 50,0.
Pode -se, também, notar um aumento do desvio padrão de 17,6, nas questões objetivas do Componente de Formação Geral dos alunos de todo o Brasil, para 21,3 nas questões discursivas do mesmo componente.
Como já comentado, a mediana de todo o Brasil, neste componente, foi 46,0, sendo esta encontrada na região Sul, enquanto a maior mediana foi encontrada na região Sudeste ( 49,5 ) e a menor na região Norte ( 42,0 ).
A maior nota máxima foi obtida, nas questões discursivas do Componente de Formação Geral, por pelo menos um aluno na região Sudeste ( 90,0 ), enquanto a menor nota máxima foi obtida na região Nordeste ( 81,5 ).
Em todas as regiões a nota mínima foi zero, sem exceção.
O Gráfico 3.15 representa a distribuição das notas nas questões discursivas no Componente de Formação Geral.
A moda desta distribuição ocorre no intervalo ( 50 ; 60 ].
Devido à grande quantidade de notas zero e à alta frequência de alunos que deixaram este tipo de questão em branco o intervalo [ 0 ; 10 ] se apresenta como uma segunda moda.
A distribuição possui assimetria à esquerda, coeficiente de assimetria -0,66.
Em todas as regiões o coeficiente de assimetria também é negativo, o maior em módulo, -0,89, na região Sudeste, e o menor na região Norte, -0,32.
A avaliação das questões discursivas de Formação Geral considerou em separado o conteúdo ( \textbf{peso} 0,8 na nota ) e o desempenho em língua portuguesa ( \textbf{peso} 0,2 ).
Na sequência, os resultados verificados para cada uma das questões discursivas de Formação Geral com relação ao conteúdo serão apresentados, estabelecendo -se relações com os temas abordados em cada uma delas.
Os comentários da Banca de docentes corretores a respeito do observado na correção das respostas dos estudantes, suas impressões e conclusões serão apresentados junto à análise de cada questão.
Em seguida será feita uma análise do desempenho de Língua Portuguesa.
Os comentários da Banca de docentes corretores serão apresentados para o conjunto de questões.
Cumpre esclarecer que, tendo em vista que as questões discursivas de Formação Geral são padronizadas, ou seja, constam de todas as provas, os comentários da Banca são os mesmos para todas as carreiras acadêmicas, sendo direcionados a todos os estudantes que participaram do ENADE/2013.
A seguir, serão analisados os desempenhos da Área de Zootecnia nas duas questões discursivas de Formação Geral do ENADE/2013, comparando os resultados obtidos com comentários para cada questão. 3.3.1.1 Análise de Conteúdo Questão Discursiva 1 do Componente de Formação Geral Os dados de Zootecnia, obtidos a partir das respostas à questão 1, encontram -se na Tabela 3.9 e no Gráfico 3.16.
Nessa questão – de melhor desempenho dentre as duas de Formação Geral – os alunos de todo Brasil tiveram média, 41,8.
A maior média para a questão 1 foi obtida na região Sudeste ( 47,1 ), e a menor, na região Norte ( 37,5 ).
Quanto à variabilidade das notas, o desvio padrão de todo o Brasil foi 25,3.
O menor desvio padrão foi obtido na região Sul ( 24,1 ) e o maior desvio padrão foi obtido na região Centro-Oeste ( 26,3 ).
A mediana da questão discursiva 1 foi 50,0 para as regiões Sudeste e Sul, e mais baixa tanto para a região Norte ( 40,0 ) quanto para as regiões Nordeste e Centro-Oeste ( 45,0 ).
Para o Brasil como um todo foi também 50,0.
As notas máximas ( 100,0 ) foram as mesmas para quase todas as regiões do Brasil, exceto as regiões Sul ( 85,0 ) e Centro-Oeste ( 90,0 ).
Além disso, a nota mínima foi zero em todas as regiões do país, sem exceção.
O Gráfico 3.16 mostra a distribuição das notas na questão discursiva 1 do Componente de Formação Geral.
Observa -se que a maior frequência corresponde aos alunos que obtiveram nota no intervalo ( 40 ; 50 ], correspondendo à moda da distribuição.
Destacam -se, também, as questões deixadas em branco e o intervalo ( 60 ; 70 ], que se apresentam como máximos locais.
A distribuição das notas possui assimetria negativa ( -0,42 ), assim como em todas as regiões do Brasil ( entre -0,61 para o Sudeste e -0,02 para o Norte ). 3.3.1.2 Comentários sobre a correção de Conteúdo das respostas à Questão Discursiva 1 O enunciado restou claro, elaborado com linguagem técnica, estando inserido nas Diretrizes gerais, com nível médio de profundidade, já que abordava tema vinculado à Área de formação da maior parte dos examinandos.
A divisão das perguntas também ficou clara indicando, passo a passo, o caminho de resposta a seguir, com base no texto.
Verificou -se dificuldade quanto à análise dos programas de governo existentes que contribuem para melhoria do saneamento básico.
Nesse ponto, os alunos referiram, com muita frequência, a ações relacionadas à Saúde, tais como : médico da família, ações educativas, ações ambientais, melhoria nas condições de habitação ( sem necessariamente vincular a programas de governo federal ou estaduais ou municipais ).
Os programas citados foram : Programa Saúde da Família ( PSF ), Programa de \textbf{Aceleração} do Crescimento ( PAC ) e Minha Casa Minha Vida e, ainda, o SUS.
Apesar de uma minoria de examinandos ter indicado o PAC do Governo Federal, a maior parte deles apresentou alguma solução para a questão do saneamento básico.
Os estudantes indicaram, com facilidade, a questão da necessidade do tratamento da água como item essencial do saneamento básico.
Muitos concluintes, através da sua experiência local, citaram situações em que as crianças caminhavam de pés no \textbf{chão} perto de rotas de esgoto a céu aberto.
Também criticaram a inércia governamental em não combater as invasões ilícitas, em locais inadequados, sem estrutura de água potável e também de saneamento básico.
Diante das \textbf{observações} preliminares pode ser extraída uma classificação das respostas obtidas que ficou estabelecida em fracas, médias e boas.
As notas consideradas fracas ( notas de 0 a 3,0 ) foram conferidas àqueles que não conseguiram relacionar o saneamento básico com a importância do tratamento de água, os projetos para resolver o problema e nem indicaram soluções para a questão.
As respostas que tiveram média avaliação ( notas de 3,5 a 7,0 ) tocaram, de algum modo, nos pontos indicados, mesmo que de forma insuficiente.
Já as respostas consideradas boas ( notas de 7,5 a 10,0 ) preencheram os requisitos do padrão de resposta, não logrando nota máxima, na maior parte dos casos, por ausência de indicação de programa governamental, mas sugerindo soluções técnicas para o problema.
As respostas mais comuns foram as consideradas medianas e que indicaram, com frequência, que o problema do saneamento básico estaria nas classes de renda mais desfavorecida e nas regiões consideradas como periferias das grandes cidades.
Também foram frequentes respostas que citaram a necessidade de tratamento de água, acompanhada da necessidade de maiores investimentos nessa área. 3.3.1.3 Análise de Conteúdo da Questão Discursiva 2 do Componente de Formação Geral A Tabela 3.10 mostra que o desempenho dos estudantes na questão 2 ( média 38,6 ) foi inferior ao obtido na questão de número 1 ( média 41,8 ).
A região Sudeste foi aquela onde a média, nessa questão, foi maior ( 41,5 ), e a de menor média foi a região Nordeste ( 35,2 ).
Quanto à variabilidade das notas, o desvio padrão de todo o Brasil foi 26,7, superior ao obtido na questão de número 1 ( 25,3 ).
O maior desvio nessa questão foi obtido na região Nordeste ( 27,7 ), enquanto o menor foi obtido na região Sul ( 25,3 ).
A mediana foi 45,0 para o Brasil como um todo, sendo encontrada também na região Sudeste.
A maior mediana foi encontrada na região Centro-Oeste ( 50,0 ), enquanto a menor foi encontrada nas regiões Norte e Nordeste ( 40,0 ).
As notas máximas ( 100,0 ) foram as mesmas em todas as regiões do Brasil, assim como as notas mínimas ( 0,0 ), exceto a nota máxima na região Norte ( 95,0 ).
O Gráfico 3.17 mostra a distribuição das notas na questão discursiva 2 do Componente de Formação Geral.
Neste gráfico, desconsiderando -se a grande quantidade de alunos que deixaram a questão 2 em branco ( um pouco menos de 20 \% ), observa -se a moda no intervalo ( 40 ; 50 ].
O coeficiente de assimetria das notas nesta questão é negativo ( -0,14 ), indicando que a distribuição é assimétrica à esquerda.
O mesmo ocorre para a maioria das regiões do Brasil ( -0,31 nas regiões Sudeste e Centro-Oeste, e -0,23 na região Sul ).
No entanto, nas regiões Norte ( 0,08 ) e Nordeste ( 0,10 ) o coeficiente de assimetria é positivo, indicando que a distribuição é assimétrica à direita. 3.3.1.4 Comentários sobre a correção de Conteúdo das respostas à Questão Discursiva 2 O tema abordado no enunciado, apesar da complexidade, é extremamente atual.
Essa atualidade facilitou a elaboração de alguma resposta, permitindo que todos os que se dispuseram a responder pudessem apresentar, pelo menos, alguns dos elementos integrantes do padrão de resposta.
A correção, seguindo o padrão de resposta, avaliou a capacidade de elaborar um texto dissertativo que abordasse o tema sob três pontos de vista : o papel da tecnologia digital na atualidade ; a garantia dos direitos do cidadão e do Estado diante do avanço tecnológico ; e o problema da segurança digital para o indivíduo, para o Estado e para as organizações ( empresas e/ou instituições nacionais ou internacionais ).
Dada a amplitude de respostas possíveis, os alunos tiveram desempenho satisfatório.
Identificamos que os concluintes, basicamente, trilharam dois caminhos : a ) a análise da influência das inovações tecnológicas, notadamente através das redes sociais, numa perspectiva essencialmente privada, daí defluindo as questões da privacidade e da segurança ; b ) uma perspectiva de segurança nacional-estatal com criticas à atuação americana de invasão de privacidade dos dados governamentais, e à inercia do governo brasileiro em reagir ao ataque, considerado, pela maioria esmagadora dos concluintes, como indevido e antiético.
Poucos foram os que uniram as duas facetas da questão : o público e o privado. 3.3.1.5 Análise de Língua Portuguesa das Questões Discursivas do Componente de Formação Geral Os dados de Zootecnia, obtidos a partir das respostas às questões discursivas do Componente de Formação Geral no que tange à Língua Portuguesa, encontram -se na Tabela 3.11 e no Gráfico 3.18.
Nesse aspecto, os alunos de todo Brasil tiveram média, 45,9.
A maior média com respeito à Língua Portuguesa foi obtida na região Sudeste ( 48,2 ), e a menor, na região Norte ( 43,0 ).
Quanto à variabilidade das notas, o desvio padrão de todo o Brasil foi 22,3.
O menor desvio padrão foi obtido na região Centro-Oeste ( 20,7 ) e o maior desvio padrão foi obtido na região Norte ( 23,6 ).
A mediana da nota de Língua Portuguesa foi 52,5 para a região Sul e Centro-Oeste, mais alta para a região Sudeste ( 55,0 ) e mais baixa para a Norte e a Nordeste ( 50,0 ).
Para o Brasil como um todo foi também 52,5.
A maior nota máxima foi obtida na região Sul ( 90,0 ), enquanto a menor nota máxima foi obtida na região Centro-Oeste ( 80,0 ).
Além disso, a nota mínima foi zero em todas as regiões do país, sem exceção. segurança ; b ) uma perspectiva de segurança nacional-estatal com criticas à atuação americana de invasão de privacidade dos dados governamentais, e à inercia do governo brasileiro em reagir ao ataque, considerado, pela maioria esmagadora dos concluintes, como indevido e antiético.
Poucos foram os que uniram as duas facetas da questão : o público e o privado.
Na maior parte das provas, o concluinte indicou a ausência de segurança e de privacidade decorrentes das próprias atividades individuais.
Muitos respondentes citaram que, de forma voluntária, os próprios cidadãos enfraquecem a proteção à privacidade, pois divulgam sua vida pessoal, aspectos de seu cotidiano, por meio de imagens e do fornecimento de dados pessoais.
No aspecto segurança, foram inúmeras as indicações de acessos indevidos a dados pessoais com a perpetração de fraudes de todos os tipos, notadamente as bancárias e comerciais, com a aquisição de bens à revelia dos titulares dos dados.
Também no tema segurança, muitos indicaram o aspecto positivo decorrente da vigilância generalizada, não somente através da rede de computadores, mas também pela quantidade de \textbf{câmeras} existentes, que contribuem para a identificação de criminosos procurados ou dos que realizam crimes pela internet.
Na perspectiva pública, a esmagadora maioria indicou repulsa pelo comportamento americano de invasão dos dados do governo brasileiro, reclamando reação do mesmo.
Alguns poucos justificaram a intervenção americana pela “ natural ” \textbf{disputa} de poder pelos Estados imperialistas. 3.3.1.5 Análise de Língua Portuguesa das Questões Discursivas do Componente de Formação Geral Os dados de Zootecnia, obtidos a partir das respostas às questões discursivas do Componente de Formação Geral no que tange à Língua Portuguesa, encontram -se na Tabela 3.11 e no Gráfico 3.18.
Nesse aspecto, os alunos de todo Brasil tiveram média, 45,9.
A maior média com respeito à Língua Portuguesa foi obtida na região Sudeste ( 48,2 ), e a menor, na região Norte ( 43,0 ).
Quanto à variabilidade das notas, o desvio padrão de todo o Brasil foi 22,3.
O menor desvio padrão foi obtido na região Centro-Oeste ( 20,7 ) e o maior desvio padrão foi obtido na região Norte ( 23,6 ).
A mediana da nota de Língua Portuguesa foi 52,5 para a região Sul e Centro-Oeste, mais alta para a região Sudeste ( 55,0 ) e mais baixa para a Norte e a Nordeste ( 50,0 ).
Para o Brasil como um todo foi também 52,5.
A maior nota máxima foi obtida na região Sul ( 90,0 ), enquanto a menor nota máxima foi obtida na região Centro-Oeste ( 80,0 ).
Além disso, a nota mínima foi zero em todas as regiões do país, sem exceção.
O Gráfico 3.18 mostra a distribuição das notas de Língua Portuguesa do Componente de Formação Geral.
Observa -se que a maior frequência corresponde aos alunos que obtiveram nota no intervalo ( 50 ; 60 ], correspondendo à moda da distribuição.
As questões em branco constituem uma moda secundária.
A distribuição das notas possui assimetria negativa ( -0,90 ), assim como em todas as regiões do Brasil ( entre -1,18 para o Sudeste e -0,64 para o Norte ). 3.3.1.6 Comentários sobre a correção das respostas de Formação Geral com respeito à Língua Portuguesa Um aspecto interessante a ser destacado é a diferença de desempenho entre a questão 1 e a questão 2, principalmente nos aspectos textuais – os textos produzidos para a questão 1 apresentaram resultados melhores do que os produzidos para a questão 2.
Uma hipótese que provavelmente explique essa discrepância é o fato de o tema “ saneamento básico ” ser mais relevante para profissionais de saúde do que o tema “ segurança e privacidade ”, constituindo tema trabalhado no currículo escolar.
Prova disso é que quase todos os textos da questão 1 citam termos específicos da área de saúde, como “ contaminação feco-oral ”, “ higiene corporal ”, “ ingestão de alimentos ou água contaminada por vetores biológicos ”.
Os resultados da avaliação correspondem aos seguintes aspectos observados em cada competência : a ) aspectos ortográficos : o desempenho dos participantes ficou concentrado, majoritariamente, nos níveis 3 ( até 3 desvios ) e 2 ( de 4 a 7 desvios ), devido ao grande índice de desvios de acentuação, contrabalançada pelo pequeno número de desvios em grafia.
Foi reduzido o número de textos com ausência de desvios ( enquadrados no nível 4 ).
Alguns textos, com pior desempenho, foram enquadrados no nível 1 ( de 8 a 12 desvios ), enquanto o nível zero foi atribuído a pouquíssimos casos, já que ele índice de desvios de acentuação, contrabalançada pelo pequeno número de desvios em grafia.
Foi reduzido o número de textos com ausência de desvios ( enquadrados no nível 4 ).
Alguns textos, com pior desempenho, foram enquadrados no nível 1 ( de 8 a 12 desvios ), enquanto o nível zero foi atribuído a pouquíssimos casos, já que ele revela falta absoluta de domínio das convenções ortográficas.
Observou -se, portanto, que existe uma diferença muito grande de desempenho nos dois aspectos analisados : baixíssimo índice de desvios ortográficos e grande índice de desvios de acentuação.
Em vários casos, ocorre ausência completa de acentuação gráfica.
Os resultados revelam, portanto, que a tendência dominante entre os universitários brasileiros é a eliminação da acentuação gráfica, talvez motivada pela ausência de esclarecimento dos meios de comunicação, das autoridades e das escolas sobre as decisões do Acordo Ortográfico da Língua Portuguesa de 1990.
Os casos mais sistemáticos de ausência de acentuação são :  palavras proparoxítonas ( “ politicas ”, “ publicas ”, “ juridicas ” ) ;  palavras paroxítonas terminadas em ditongo crescente ( “ varios ”, “ Brasilia ”, “ familia ” ) ;  palavras oxítonas ( “ ninguem ”, “ esta ”, “ ate ” ).
Quanto ao domínio das convenções relativas à grafia das palavras, são pontuais os problemas encontrados ( “ infecciozas ”, “ obtos ”, “ hipedemia ”, “ saniamento ”, “ cituação ”, “ abitantes ” ).
Os únicos desvios recorrentes foram : a ausência de segmentação nos casos de verbo + pronome oblíquo “ se ”, combinada, eventualmente, à alteração da grafia do pronome ( “ percebese ”, “ percebece ”, “ preferece ” ) ; a grafia da forma verbal de pretérito imperfeito do subjuntivo como se fosse uma ocorrência de verbo + pronome oblíquo ( “ distribui -se ” por “ distribuísse ” ) Vale observar, também, que, ao contrário do que se esperava, não apareceram abreviaturas próprias do “ internetês ”, ou seja, dos hábitos de comunicação escrita adquiridos pelo uso de redes sociais e e-mails. b ) aspectos textuais : o desempenho dos participantes ficou concentrado, majoritariamente, nos níveis 3 e 2, devido à grande ocorrência de problemas de estruturação textual.
Foi muito baixo o número de textos enquadrados no nível 4 por não apresentar qualquer problema estrutural.
Alguns textos com pior desempenho foram enquadrados no nível 1, enquanto o nível zero foi atribuído a pouquíssimos casos, já que ele revela a existência de um texto sem articulação e com comprometimento do sentido.
Observou -se que a grande maioria dos participantes não distribuiu as ideias em parágrafos, talvez devido ao pequeno número de linhas disponibilizadas para a resposta da questão ou, quem sabe, pela suposição de que não seria necessária essa divisão por não se tratar de um texto no modelo de uma redação dissertativo-argumentativa, como solicitado nos vestibulares.
Essa competência é a que se revela como a mais problemática entre os participantes, porque são muitos os problemas observados, desvios acumulados durante toda a formação do estudante e que não se resolvem com um estudo autodidata, como acontece com regras ortográficas ou morfossintáticas : sequência justaposta de ideias sem encaixamentos sintáticos ; redução drástica de estruturas subordinadas, ao lado do aumento na frequência de estruturas coordenadas e absolutas ; redução no uso de conectores para expressar relações lógicas essenciais à construção do texto, substituídas pela exigência de inferência por parte do interlocutor para suprir a sua ausência ; emprego equivocado de operadores que não estabelecem relações lógicas coerentes entre ideias do texto ; emprego inadequado do pronome relativo ( com omissão da preposição ou a utilização de pronome inadequado, como “ onde ” ) ; repetição exaustiva de termos sem a utilização de procedimentos mais sofisticados de substituição ( hiperonímias, hiponímias, nominalizações, expressões metafóricas ) ; frases fragmentadas que comprometem a estrutura lógico-gramatical ; frases formadas apenas por oração subordinada, sem oração principal.
Um importante aspecto a destacar é o baixíssimo desempenho de uma parte dos participantes em relação à estrutura formal do texto produzido, o que é extremamente preocupante ao se levar em conta que são graduandos em fase final de formação.
São frequentes os casos de desvios de estruturação frasal, com uso inadequado ou ausência de conectivos entre parágrafos e entre frases.
Em uma parte dos textos, falta um mínimo de textualidade e de domínio do registro padrão da língua.
Na verdade, observam -se relações linguísticas quase agramaticais, como as estabelecidas pela sequência de gerúndios sem o apoio de um ponto de partida para a organização das informações gramaticais e semânticas.
Quanto à utilização dos mecanismos de referenciação, dois aspectos devem ser destacados : a ocorrência, em uma boa parte dos textos, de repetições de palavras ou expressões sem a utilização de termos sinônimos ou pronomes, como seria adequado ; e a alternância entre os pronomes demonstrativos este / esse, e suas formas flexionadas, sem a observância de um critério de utilização.
Esse último aspecto tem se generalizado na modalidade oral da língua e, cada vez mais, tem se tornado uma constante na modalidade escrita de jornais e revistas : predominância da forma “ esse ” na modalidade oral e da forma “ este ” na modalidade escrita, provavelmente por um processo de hipercorreção. c ) aspectos morfossintáticos e vocabulares : o desempenho dos participantes ficou concentrado, majoritariamente, nos níveis 3 e 2, devido à grande ocorrência de problemas de regência, concordância e uso adequado dos sinais de pontuação.
O nível 4 foi atribuído a um número muito reduzido de textos.
Alguns textos com pior desempenho foram enquadrados no nível 1, enquanto o nível zero foi atribuído a pouquíssimos casos, já que ele revela a existência de um texto sem articulação e com comprometimento do sentido.
Os resultados são muito transparentes em relação aos aspectos mais problemáticos no desempenho dos participantes.
O desvio mais frequente é a falta do sinal indicativo da crase, o que se configura como um problema de regência, verbal ou nominal, conforme o caso – isso revela que o usuário não tem consciência de que, sob a forma do termo “ a ” existe a presença de uma preposição “ a ”, exigida pela regência do termo anterior.
Embora em outros exames, como o Enem, a falta de crase seja \textbf{apenada} em acentuação, nesta correção amostral esse desvio foi considerado no âmbito dos aspectos morfossintáticos.
Outro problema relacionado à regência verbal e à nominal, encontrado frequentemente nas questões, foi a ausência de preposição antes de pronome relativo, processo generalizado na modalidade oral da língua, em situações de registro informal.
Apesar da possibilidade de que essa alteração de regência se generalize no padrão escrito da Língua Portuguesa, como já está ocorrendo até em textos jornalísticos, o não emprego da preposição foi \textbf{apenado} neste processo de avaliação.
Outro desvio muito frequente diz respeito aos processos de concordância verbal e de concordância nominal.
Quanto à concordância de número, observou -se ausência de marca ( com sujeito anteposto ou posposto ) ou uso indevido ( uso inadequado da marca de plural comandado pelo núcleo plural da locução adjetiva, apesar de o substantivo que funciona como núcleo do sintagma nominal estar no singular ).
Uma ocorrência que se destacou foi a ausência de acento circunflexo na forma plural do presente do indicativo do verbo “ ter ”, que foi considerada como um desvio na concordância verbal e não na acentuação gráfica.
Quanto à concordância de gênero, vários casos foram observados, normalmente no âmbito de sintagmas nominais longos, em que o adjetivo está afastado do substantivo.
Quanto à questão da colocação pronominal, foram poucos os casos observados.
Apesar de serem aspectos relacionados à oralidade, concluiu -se que, no registro escrito formal, a maioria dos participantes já incorporou regras como a não introdução da frase por um pronome oblíquo e a próclise na presença de um termo atrator.
Não se adotou, entretanto, o padrão excessivamente formal descrito pelas gramáticas normativas em relação à posição do pronome oblíquo em locuções verbais, já que esse uso está muito distante da prática cotidiana, até em textos mais formais.
Quanto aos aspectos vocabulares, alguns tipos de inadequação foram observados : expressões da oralidade apareceram em algumas respostas, mas sem maior relevância do ponto de vista quantitativo ; seleção vocabular incompatível com o contexto, gerando situações de falta de inteligibilidade ; falta de domínio de vocabulário mais abstrato e de maior complexidade, essencial ao desenvolvimento do texto de base dissertativa.
Quanto à utilização dos sinais de pontuação, observou -se uma grande precariedade nos textos analisados.
Não foi \textbf{apenada} a ausência de vírgula para destacar locuções ou adjuntos adverbiais de pequena extensão deslocados de posição na frase, por ser um uso opcional.
São os seguintes os tipos de problemas encontrados : a ) vírgula : utilização de vírgula para separar o sujeito e o predicado ; ocorrência de apenas uma das vírgulas para separar uma palavra, uma expressão ou uma oração encaixada ; uso de vírgula no lugar do ponto para separar ideias que constituem períodos distintos ; ausência de vírgula para separar enumerações ; ausência de vírgula para separar oração adjetiva explicativa ou utilização inadequada para separar oração adjetiva restritiva ; b ) ponto e vírgula : utilização do ponto e vírgula no lugar de vírgula ; c ) ponto final : ausência de ponto final para separar períodos.
Várias marcas de oralidade foram identificadas, embora não em alta frequência : o uso do pronome relativo “ onde ” como relativo universal, falta de artigo definido antes de substantivo, repetição de palavras por falta de vocabulário, reduções como “ tá ”, “ pra ”, “ pro ”, “ prum ”, expressões informais como “ estão nem aí ”, “ tem que ”, “ levando com a barriga ”, “ não tá nem aí ”, “ a coisa fica séria ”, “ e aí haja dor de cabeça ”. 3.3.2 Componente de Conhecimento Específico Na parte da prova relativa às questões discursivas no Componente de Conhecimento Específico ( Tabela 3.12 ), observa -se que a média foi mais baixa do que nas questões discursivas do Componente de Formação Geral.
Enquanto no Componente de Formação Geral a média para estudantes de Zootecnia de todo o Brasil foi 41,3, na Parte de Conhecimento Específico a média foi 30,0.
A maior média deste componente foi obtida pelos estudantes da região Sudeste ( 35,4 ), e a menor, pelos da região Norte ( 21,9 ).
Quanto à variabilidade das notas, o desvio padrão de todo o Brasil foi 19,6.
O maior desvio padrão foi encontrado nas regiões Sul e Centro-Oeste ( 19,8 ), e o menor, na região Nordeste ( 17,8 ).
A maior nota máxima, nas questões discursivas do Componente de Conhecimento Específico, foi obtida na região Sudeste ( 100,0 ), enquanto a menor nota máxima foi encontrada na região Norte ( 85,0 ).
Além disso, a nota mínima ( 0,0 ) foi obtida por alunos de todas as regiões do Brasil sem exceção.
A mediana do Brasil como um todo foi 30,0, também encontrada na região Centro-Oeste.
As demais medianas foram : 20,0 na região Norte, 23,3 ( a mais baixa ) na região Nordeste e 35,0 nas regiões Sudeste e Sul ( as mais altas ).
O Gráfico 3.19 representa a distribuição das notas nas questões discursivas no Componente de Conhecimento Específico.
Essa distribuição tem moda no intervalo de ( 20 ; 30 ].
Nota -se que há uma tendência decrescente a partir deste intervalo até o último intervalo.
O intervalo [ 0 ; 10 ] se apresenta como uma moda secundária.
O coeficiente de assimetria é positivo para todo o Brasil ( 0,30 ) e para todas as regiões ( entre 0,11 para a região Sul e 0,67 para a região Norte ), indicando uma assimetria à direita. 3.3.2.1 Análise da Questão Discursiva 3 do Componente de Conhecimento Específico Na questão 3, cujos resultados aferidos encontram -se descritos na Tabela 3.13, a média dos estudantes de todo o Brasil foi 14,3.
O desempenho médio dos estudantes nesta questão foi inferior ao desempenho médio nas questões de número 4 e 5.
A menor média nessa questão foi obtida pelos alunos da região Nordeste ( 7,6 ), enquanto a maior média foi obtida na região Sudeste ( 18,9 ).
Quanto à variabilidade das notas, o desvio padrão do Brasil foi 23,5.
O maior desvio padrão foi obtido na região Centro-Oeste ( 26,2 ), enquanto o menor foi obtido na região Nordeste ( 18,0 ).
A nota máxima, 100,0 pontos, foi alcançada por pelo menos um aluno em todas as regiões, sem exceção.
A mediana e a nota mínima do Brasil como um todo foi zero, o mesmo valor para todas as regiões, exceto a mediana nas regiões Sudeste e Sul ( 5,0 ).
Ou seja, exceto nas regiões Sudeste e Sul, pelo menos metade dos estudantes de cada região recebeu nota zero nesta questão, seja por terem deixado a questão em branco, seja pela inadequação da resposta.
O Gráfico 3.20 mostra a distribuição das notas na questão discursiva 3, do Componente de Conhecimento Específico.
Essa distribuição é unimodal com moda nas questões em branco.
O coeficiente de assimetria é positivo para todo o Brasil ( 1,85 ) e para todas as regiões ( entre 1,38 na região Sudeste e 2,85 na região Nordeste ), indicando uma assimetria acentuada à direita. 3.3.2.2 Comentários sobre as respostas à Questão Discursiva 3 O enunciado foi claro, de modo que os estudantes não demonstraram dúvidas quanto ao que foi solicitado na questão.
A profundidade do conhecimento exigido foi adequada, sendo a questão de média-alta complexidade, uma vez que tratava de um ponto específico - e muito importante - da Zootecnia que envolvia cálculos matemáticos.
Desse modo, a questão tinha como resposta valores exatos, apesar de alguma margem para aproximação.
Além disso, os estudantes que desconheciam o caminho para obtenção da resposta simplesmente deixaram a questão em branco.
Não foram observadas divergências significativas entre as respostas dos estudantes e o padrão de resposta ; provavelmente pelo fato de a questão envolver cálculos matemáticos e possuir como respostas valores numéricos exatos.
As pequenas variações observadas foram relacionadas a arredondamentos de decimais e estavam previstas no padrão de resposta.
A distribuição da pontuação foi acertada, com subitens mais complexos recebendo maior pontuação que os subitens mais simples.
De modo geral o desempenho dos alunos foi fraco.
Foram frequentemente encontradas respostas com textos gerais, sem mostrar conhecimento sobre o tema. 3.3.2.3 Análise da Questão Discursiva 4 do Componente de Conhecimento Específico A Tabela 3.14 contém as informações relativas à questão 4 do conjunto de questões do Componente de Conhecimento Específico.
O desempenho médio dos estudantes nesta questão foi superior ao desempenho nas questões de número 3 e 5.
A média geral do Brasil foi 38,9, sendo a menor média registrada na região Norte ( 31,6 ) e a maior na região Sudeste ( 44,4 ).
Quanto à variabilidade das notas, o desvio padrão de todo o Brasil foi 25,5.
O maior desvio padrão foi obtido na região Centro-Oeste ( 26,3 ), enquanto o menor foi obtido na região Nordeste ( 24,3 ).
A nota máxima ( 100,0 ) foi atingida em quase todas as regiões, exceto nas regiões Norte ( 90,0 ) e Nordeste ( 95,0 ).
A mediana do Brasil como um todo foi 45,0, sendo encontrada também na região Centro-Oeste.
A maior mediana foi encontrada nas regiões Sudeste e Sul ( 50,0 ), enquanto a menor mediana foi encontrada nas regiões Norte e Nordeste ( 35,0 ).
A nota mínima foi zero para todas as regiões do Brasil, sem exceção.
O Gráfico 3.21 representa a distribuição de notas na questão discursiva 4, no Componente de Conhecimento Específico.
Como nas outras questões discursivas, o número de estudantes que deixou a questão em branco foi muito grande ( em torno de 18 \% ).
A moda desta distribuição ocorre no intervalo ( 40 ; 50 ].
O coeficiente de assimetria é negativo para todo o Brasil ( -0,09 ) e para quase todas as regiões ( -0,13 para a região Sudeste, -0,27 para a região Sul e -0,16 para a região Centro-Oeste ), indicando uma assimetria à esquerda.
Já as regiões Norte ( 0,00 ) e Nordeste ( 0,04 ) possuem assimetria nula ou positiva. 3.3.2.4 Comentários sobre as respostas à Questão Discursiva 4 A Questão 4 abordava um tema de grande relevância para a área de Zootecnia : a importância e os desafios do Agronegócio.
O enunciado era de fácil entendimento, ao mesmo tempo em que exigia um nível adequado de profundidade de conhecimentos.
O enunciado fornecia um grande número de informações, não deixando margem a interpretações diferentes.
No entanto, as informações - em excesso - foram utilizadas pelos estudantes na elaboração de respostas que nada acrescentam ao que já havia sido exposto.
Nesse sentido, na primeira parte da questão, item ‘ a ’, grande parte dos estudantes apenas repetiu as informações contidas no enunciado - citando, por exemplo, o termo “ PIB ” de forma isolada e descontextualizada.
O padrão de respostas utilizado para a correção era abrangente e adequado à questão, sendo a pontuação bem distribuída entre os tópicos.
Sendo assim, as respostas dos estudantes, de modo geral, estavam previstas no padrão de respostas.
No entanto, o tema “ Sustentabilidade ” foi muito citado pelos estudantes, mas não foi previsto explicitamente no padrão de respostas.
Na maior parte das respostas foi dada muita ênfase ao aspecto econômico do Agronegócio, em detrimento do aspecto social.
Possivelmente, esse foi mais um reflexo do grande número de dados econômicos apresentados no enunciado.
As respostas mais frequentemente encontradas foram aquelas com textos gerais, sem aprofundamento suficiente para mostrar conhecimento sobre o tema.
Foi possível observar que uma pequena parcela dos estudantes recorreu a questões objetivas para elaborar sua resposta, particularmente a alternativa ‘ a ’.
Dentre essas questões objetivas, podemos citar : a de nº 02, que abordou a temática “ desenvolvimento em bases sustentáveis ” ; a de nº 07 que abordou a temática “ projeto de uma sociedade sustentável ” ; e por fim a de nº 22 que abordou a temática “ A sustentabilidade ambiental e econômica de pequenas propriedades rurais ”. 3.3.2.5 Análise da Questão Discursiva 5 do Componente de Conhecimento Específico A Tabela 3.15 contém as informações relativas à questão 5 do conjunto do Componente de Conhecimento Específico.
O desempenho dos estudantes nessa questão foi superior ao da questão 3 e inferior ao da questão 4.
A nota média dos estudantes de todo o Brasil foi 36,9.
A maior média foi registrada na região Sul ( 43,1 ), enquanto a menor média foi registrada na região Norte ( 24,6 ).
Quanto à variabilidade das notas, o desvio padrão dos alunos do Brasil, como um todo, foi 30,9.
Enquanto o maior desvio foi encontrado na região Sudeste ( 32,2 ), o menor foi encontrado na região Norte ( 26,3 ).
A nota máxima ( 100,0 ) foi alcançada em todas as regiões do Brasil, sem exceção.
A mediana foi 40,0 para quase todas as regiões do Brasil, exceto para as regiões Norte ( 20,0 ) e Nordeste ( 25,0 ).
A nota mínima foi zero para todas as regiões do Brasil, sem exceção.
O Gráfico 3.22 representa a distribuição das notas da questão discursiva 5 do Componente de Conhecimento Específico.
Destaca -se, como nas demais questões discursivas, o grande número de estudantes que deixaram esta questão em branco, correspondendo à moda da distribuição.
Além da moda na questão em branco, aparecem ainda três máximos locais com representação acima de 10,0 \% do total avaliado : ( 10 ; 20 ], ( 30 ; 40 ] e ( 50 ; 60 ].
O coeficiente de assimetria para todos os alunos é 0,40 e para todas as regiões se mantém positivo, variando entre 0,13 na região Sul e 0,90 na região Norte. 3.3.2.6 Comentários sobre as respostas à Questão Discursiva 5 A questão versava sobre tópico que frequentemente volta a ser abordado nos meios de comunicação.
Para elaborar uma resposta correta não era necessário fazer associações complexas.
A resposta era simples e exigia apenas capacidade de resgatar conhecimentos fundamentais do curso de Zootecnia.
Na verdade, o item a poderia ser respondido por qualquer profissional da área agrária, bastando que tivesse acesso a textos técnicos como o veiculado na página da UBABEF, em 14/11/2013, sobre o tema.
O padrão de resposta do subitem a contempla o que poderia ser considerada a resposta ideal, adequando -se ao que foi perguntado.
Neste subitem, houve pouca variação nas respostas.
A maioria das respostas do subitem ‘ a ’ foi muito superficial e não abordou o tema na profundidade exigida.
Quanto ao subitem ‘ b ’, o enunciado da pergunta pedia para que fossem citados os aditivos que são permitidos pela legislação brasileira, o que causou bastante confusão nas respostas.
Nesse subitem, grande parte dos estudantes não citou corretamente todas as classes de aditivos zootécnicos, assim como não explicitou suas funções.
A subdivisão da pontuação foi adequada de acordo à complexidade dos subitens.
Observou -se que um número considerável de alunos identificou -se ao responder a questão e, portanto, tiveram suas respostas desconsideradas.
Esse fato talvez tenha sido motivado pela formulação do item ‘ a ’ que solicitava “ elabore uma nota de esclarecimento à população ” redigir um texto identificado.

% matched lemmas: aceleração, apenar, chão, câmera, disputar, observação, peso
\end{textsample}
